%! Matrices and Column Spaces — Unit 1, Lesson 1.3 — Guided Notes
\documentclass[10pt]{article}
\usepackage{linalg-article}
\usepackage{linalg-boxes}
\newcommand{\TallMath}[1]{$\displaystyle #1\rule[-1.4em]{0pt}{3.2em}$}
\newcommand{\vv}{\mathbf{v}}
\newcommand{\ww}{\mathbf{w}}
\newcommand{\xx}{\mathbf{x}}
\newcommand{\bb}{\mathbf{b}}

\begin{document}

\pageheader{Unit 1, Lesson 1.3}{Guided Notes}
\namedateperiod

\vspace{0.12in}

\begin{objectivebox}
\textbf{By the end of this lesson, I will be able to\dots}
\begin{itemize}\setlength{\itemsep}{2pt}
  \item read a \emph{matrix} as its \emph{columns} and compute $A\xx$ as a combination of them;
  \item describe the \emph{column space} $C(A)$ and tell whether it is a line or the whole plane;
  \item decide whether a target $\bb$ is \emph{reachable} (in the column space).
\end{itemize}
\end{objectivebox}

\vspace{0.10in}

\begin{vocabbox}
\small
Fill in each term as we build it together.
\vspace{2pt}
\termblanklong{Matrix (its columns)}
\termblanklong{Matrix--vector product $A\xx$}
\termblanklong{Column space $C(A)$}
\termblanklong{Independent / dependent columns}
\end{vocabbox}

\begin{hookbox}
A matrix is really just two vectors standing side by side:
\[
  A=\begin{bmatrix} 1 & 3 \\ 2 & 1 \end{bmatrix}
  \quad\text{has columns}\quad
  \mathbf{a}_1=\begin{bmatrix} 1 \\ 2 \end{bmatrix},\ \
  \mathbf{a}_2=\begin{bmatrix} 3 \\ 1 \end{bmatrix}.
\]
Multiplying by $\xx=\begin{bsmallmatrix} 2\\1 \end{bsmallmatrix}$ \emph{combines} the columns:
$2\begin{bsmallmatrix} 1\\2 \end{bsmallmatrix}+1\begin{bsmallmatrix} 3\\1 \end{bsmallmatrix}=\blank{2.0cm}$.
As $\xx$ changes, which targets can we build? \writeline
\end{hookbox}

\begin{notesbox}{1. A Matrix Is Vectors in Columns}
A \textbf{matrix} is a block of numbers --- but the useful way to read it is as \emph{vectors
standing side by side}. Its \textbf{columns} are those vectors:
\[
  A=\begin{bmatrix} 1 & 3 \\ 2 & 1 \end{bmatrix}
  =\begin{bmatrix} \mathbf{a}_1 & \mathbf{a}_2 \end{bmatrix},
  \qquad
  \mathbf{a}_1=\begin{bmatrix} 1 \\ 2 \end{bmatrix},\ \
  \mathbf{a}_2=\begin{bmatrix} 3 \\ 1 \end{bmatrix}.
\]
Your turn: the columns of $\begin{bmatrix} 4 & 0 \\ 1 & 5 \end{bmatrix}$ are
$\begin{bmatrix}\rule{0pt}{1.1em}\blank{0.8cm}\\[2pt]\blank{0.8cm}\end{bmatrix}$ and
$\begin{bmatrix}\rule{0pt}{1.1em}\blank{0.8cm}\\[2pt]\blank{0.8cm}\end{bmatrix}$.
\end{notesbox}

\begin{notesbox}{2. Multiplying a Matrix by a Vector Combines the Columns}
To compute $A\xx$, use the weights in $\xx$ to scale the columns, then add:
\[
  A\xx = \begin{bmatrix} \mathbf{a}_1 & \mathbf{a}_2 \end{bmatrix}
  \begin{bmatrix} x_1 \\ x_2 \end{bmatrix}
  = x_1\,\mathbf{a}_1 + x_2\,\mathbf{a}_2 .
\]
This is a \textbf{linear combination of the columns} --- exactly the Lesson~1.1 idea.

\vspace{2pt}
\textit{Example.} $A=\begin{bmatrix} 1 & 3 \\ 2 & 1 \end{bmatrix}$, $\xx=\begin{bmatrix} 2 \\ 1 \end{bmatrix}$:
\[
  A\xx = 2\begin{bmatrix} 1 \\ 2 \end{bmatrix} + 1\begin{bmatrix} 3 \\ 1 \end{bmatrix}
  = \blank{2.4cm}.
\]

\vspace{2pt}
Your turn: $\begin{bmatrix} 4 & 0 \\ 1 & 5 \end{bmatrix}\begin{bmatrix} 1 \\ 2 \end{bmatrix}
= 1\begin{bmatrix} 4 \\ 1 \end{bmatrix} + 2\begin{bmatrix} 0 \\ 5 \end{bmatrix} = \blank{2.0cm}$.
\end{notesbox}

\begin{notesbox}{3. The Column Space: Everything the Columns Can Build}
Let $\xx$ be \emph{anything}. The set of all outputs $A\xx$ is the \textbf{column space}
$C(A)$ --- every linear combination of the columns, i.e.\ their \textbf{span}.
\begin{center}
\begin{tikzpicture}[scale=0.62]
  \fill[burgundy!10] (-2.3,-2.3) rectangle (2.3,2.3);
  \draw[step=1, linegray!60, thin] (-2.3,-2.3) grid (2.3,2.3);
  \draw[->, thick, charcoal] (-2.5,0)--(2.5,0) node[right]{\scriptsize $x$};
  \draw[->, thick, charcoal] (0,-2.5)--(0,2.5) node[above]{\scriptsize $y$};
  \draw[->, very thick, burgundy] (0,0)--(1,2) node[above]{\scriptsize $\mathbf{a}_1$};
  \draw[->, very thick, royalblue] (0,0)--(2,0.667) node[right]{\scriptsize $\mathbf{a}_2$};
  \node[charcoal] at (0,-3.05) {\scriptsize $C(A)=$ whole plane (columns not parallel)};
\end{tikzpicture}
\end{center}
Because $\mathbf{a}_1=\begin{bsmallmatrix} 1\\2 \end{bsmallmatrix}$ and
$\mathbf{a}_2=\begin{bsmallmatrix} 3\\1 \end{bsmallmatrix}$ are \emph{not parallel}, their
combinations reach \emph{every} point: $C(A)$ is the \blank{2.6cm}.
\end{notesbox}

\begin{notesbox}{4. Line vs.\ Plane --- and Is $\bb$ Reachable?}
When the two columns \emph{are} parallel, they can only build a \textbf{line}. Take
$B=\begin{bmatrix} 1 & 2 \\ 2 & 4 \end{bmatrix}$: column~2 is
$2\times$ column~1, so every $B\xx$ lands on the line through $\begin{bsmallmatrix} 1\\2 \end{bsmallmatrix}$.
\begin{center}
\begin{tikzpicture}[scale=0.62]
  \draw[step=1, linegray!60, thin] (-2.3,-2.3) grid (2.3,2.3);
  \draw[->, thick, charcoal] (-2.5,0)--(2.5,0) node[right]{\scriptsize $x$};
  \draw[->, thick, charcoal] (0,-2.5)--(0,2.5) node[above]{\scriptsize $y$};
  \draw[burgundy!55, very thick] (-1.15,-2.3)--(1.15,2.3);
  \draw[->, very thick, burgundy] (0,0)--(1,2) node[left=1pt]{\scriptsize $\begin{bsmallmatrix} 1\\2 \end{bsmallmatrix}$};
  \node[charcoal] at (0,-3.05) {\scriptsize $C(B)=$ a line (columns parallel)};
\end{tikzpicture}
\end{center}
Solving $A\xx=\bb$ asks: \emph{does $\bb$ lie in the column space?}
\begin{itemize}\setlength{\itemsep}{1pt}
  \item Is $\begin{bmatrix} 3 \\ 6 \end{bmatrix}$ in $C(B)$? \blank{2.6cm} (on the line? \blank{1.2cm})
  \item Is $\begin{bmatrix} 1 \\ 0 \end{bmatrix}$ in $C(B)$? \blank{2.6cm} (on the line? \blank{1.2cm})
\end{itemize}
Your turn: For $A=\begin{bsmallmatrix} 1&3\\2&1 \end{bsmallmatrix}$ (columns not parallel), is
every $\bb$ reachable? \writeline
\end{notesbox}

\begin{practicebox}
Let $A = \begin{bmatrix} 1 & 2 \\ 3 & 1 \end{bmatrix}$ with columns
$\mathbf{a}_1=\begin{bsmallmatrix} 1\\3 \end{bsmallmatrix}$, $\mathbf{a}_2=\begin{bsmallmatrix} 2\\1 \end{bsmallmatrix}$.
\begin{enumerate}[leftmargin=1.6em]
  \item Compute $A\begin{bmatrix} 1 \\ 2 \end{bmatrix}$ as a combination of the columns. \hfill \blank{4.0cm}
  \item Is the column space $C(A)$ a line or the whole plane? How do you know? \writeline
  \item Is $\begin{bmatrix} 2 \\ 4 \end{bmatrix}$ in the column space of $\begin{bmatrix} 1 & 3 \\ 2 & 6 \end{bmatrix}$?
        Explain. \writeline
\end{enumerate}
\end{practicebox}

\end{document}
